\subsubsection{Carving the interior balls}
\label{ha:carving}

Write $[z]_+=\max\{z,0\}$ and set
\[
 \theta=\frac{C_4}{6C_3}.
\]
The assumptions on the constants imply
\begin{equation}
 \begin{split}
 C_3-1&\geq 2\bigl(C_4+4(C_4+2)^2\bigr),\\
 \frac78(C_3-1)&\geq C_4+4,\qquad
 C_3\geq 2C_4(C_4+2),\qquad
 \theta\leq\frac1{12(C_4+2)}\leq\frac1{60}.
 \end{split}
 \label{ha:constant-consequences}
\end{equation}
For example, the first inequality follows by subtracting $1$ from
$C_3\geq2(C_4+1)+8(C_4+2)^2$; this same lower bound gives the remaining
inequalities for $C_4\geq3$.

During phase $\ell$, assign a nonnegative potential to each record
$A$ in the current arrangement, with center $x$ and numerical inner
and outer radii $a<b$:
\begin{equation}
 \operatorname{tok}_\ell(A)=
 \begin{cases}
  \delta_\ell-C_4\delta_{\ell+1},
   &A\text{ is an original ball of the current family }\mathcal F_\ell,\\
  [b-a-C_4\delta_\ell]_+,
   &A\text{ is nonavoiding and is not such an original ball},\\
  [b-a-C_4\delta_\ell]_+/\bigl(2(C_4+2)\bigr),
   &A\text{ is avoiding}.
 \end{cases}
 \label{ha:potential}
\end{equation}
Here and below, membership in a family refers to one of its original
balls, with its original record, rather than to a newly created shell
that happens to have the same support.  The first line satisfies
\begin{equation}
 \delta_\ell-C_4\delta_{\ell+1}\geq\frac56\delta_\ell.
 \label{ha:fresh-potential}
\end{equation}
Every potential is at most the width of its record.  An empty shell,
including a proposed shell whose outer radius is at most its inner
radius, is discarded and contributes zero.

We describe the processing of a fixed bag $\operatorname{Bag}_t$ in
phase $\ell$.  Put $\delta=\delta_\ell$.  Initially partition its balls
into the interior balls $\operatorname{Int}_t$ and the border balls
$\operatorname{Bor}_t$.  The sets contain only balls still awaiting
treatment.  Repeatedly perform one of the three operations below while
\begin{equation}
 |\operatorname{Int}_t|\geq\frac12|\operatorname{Bag}_t|.
 \label{ha:interior-loop}
\end{equation}
After an operation, remove the balls explicitly treated by it from
$\operatorname{Int}_t$, and transfer every remaining ball that has
ceased to be interior to $\operatorname{Bor}_t$.  Lemma~\ref{ha:border-monotone}
ensures that this transfer is permanent.  A ball removed in one of the
three operations is counted as treated even in the third operation,
where a pair of new shells pays for it without inserting that ball.

The subscript $s$ denotes the state after $s$ operations for this bag;
the subscript $\infty$ denotes the final state.  In particular, set
\[
 n_{t,s}=|\operatorname{Bag}_t|
          -|\operatorname{Bor}_{t,s}|-|\operatorname{Int}_{t,s}|.
\]
Moving a ball from the interior set to the border set does not change
$n_{t,s}$.  Treating $m$ balls increases it by exactly $m$.

\begin{lemma}[Potential preserved during interior processing]
\label{ha:construction}
Suppose phase $\ell$ starts with a valid arrangement whose potentials
can pay each earlier node $v\in T_i$, $i<\ell$, at least
$(\theta/2)|\operatorname{Bag}_v|\delta_i$.  Interior processing preserves
validity and these allocations.  At every intermediate state it also
pays each completed node $v\in T_\ell$ at least
$\theta n_{v,\infty}\delta_\ell$, and the node currently being processed
at least $\theta n_{t,s}\delta_\ell$.
Each bag finishes with
$|\operatorname{Int}_{t,\infty}|<|\operatorname{Bag}_t|/2$.
\end{lemma}

\begin{proof}
We maintain actual allocations of the potentials, so that an operation
must preserve all amounts already allocated.  If an old record $B$ is
removed, it suffices to set aside at least
$\operatorname{tok}_\ell(B)$ from its replacements: those funds can
then be distributed exactly as the old potential was.  Any additional
funds assigned to $t$ increase its allocation.  Untouched records retain
their allocations.

Shrinking a record about its existing center and retaining its labels
preserves all its individual validity conditions.  Indeed, its old
inner witness remains inside the new inner radius, its old crossing
paths still cross the smaller radial interval, and all excluded
boundary endpoints remain beyond the new outer radius.  Pairwise
validity is inherited by Lemma~\ref{ha:heredity}.  These observations
apply to every residual shell below.  Original current-family balls
are never carved: their width is $\delta$, whereas a record containing
an interior ball must have width greater than $2C_3\delta$.

\paragraph{\textbf{Operation 1: carving an avoiding shell.}}
If some $A=B(X_u,\delta)\in\operatorname{Int}_t$ lies in an avoiding
shell $B=B(x,r_1\Vert r_2)$, put $q=\|X_u-x\|_1$ and replace $B$ by
\[
 B_1=B(x,r_1\Vert q-\delta),\qquad
 B_2=B(x,q+\delta\Vert r_2),
\]
both with the labels of $B$.  Insert the original ball $A$, with its
original labels, and treat it.  The two residual shells are disjoint
from $A$.  The original ball satisfies its individual validity
conditions, and Lemma~\ref{ha:unique-interior} shows that its labels
are compatible with every other existing record.  This proves validity
of the new arrangement.

Write $R=r_2-r_1$.  Since $R>2C_3\delta$, the old potential is
$(R-C_4\delta)/(2(C_4+2))$.  Give $B$ all the potentials of $B_1,B_2$ and
$\delta/2$ from $A$.  Positive parts can only increase the following
lower bound:
\begin{align*}
 \operatorname{tok}_\ell(B_1)+\operatorname{tok}_\ell(B_2)
       +\frac\delta2
 &\geq\frac{R-2\delta-2C_4\delta}{2(C_4+2)}
       +\frac\delta2\\
 &=\frac{R-C_4\delta}{2(C_4+2)}
  =\operatorname{tok}_\ell(B).
\end{align*}
The remaining potential of $A$ is at least
\[
 \delta-C_4\delta_{\ell+1}-\frac\delta2
 \geq\frac\delta3\geq\theta\delta.
\]
Give it to $t$.  Thus treating this one ball preserves every old
allocation and increases the allocation to $t$ by the required amount.

\paragraph{\textbf{The common setup for operations 2 and 3.}}
Suppose no ball of $\operatorname{Int}_t$ is interior to an avoiding
shell.  Write $t^*$ for the parent of $t$.
If enough interior centers lie near a short radial interval, inserting
their balls will pay for the width removed from the old record.
When the interval is long compared with their total radius, we instead
create two overlapping central shells with crossing paths in
$G[V(t)]$ and $G[V(t^*)\setminus V(t)]$, respectively.
Local connectivity supplies the first set of paths; the decomposition
below supplies the second.  Define
\[
 H=G[V(t^*)\setminus V(t)],\qquad
 O'(t)=\{\{u,v\}\in O(t):\|X_u-X_v\|_1<\delta\}.
\]
Partition $V(H)$ into sets $S_1,\ldots,S_j$ whose induced graphs are
$k/4$-edge-connected, by repeatedly splitting a part along a cut of
size less than $k/4$ until no such cut remains.  Singleton parts are
allowed.  There are $j-1$ splits, and each edge between final parts is
counted at precisely the split that first separates its endpoints.
Consequently
\[
 \sum_{i=1}^j|\partial_H S_i|
    <\frac{k}{2}(j-1)\quad(j>1),
\]
and the sum is zero when $j=1$.  Each $S_i$ is a nonempty proper subset
of $V(t^*)$, so the $k$-edge-connectivity of $G[V(t^*)]$ gives
\begin{align*}
 jk
 &\leq\sum_{i=1}^j|\partial_{G[V(t^*)]}S_i|\\
 &=|O(t)|+\sum_{i=1}^j|\partial_H S_i|
 \leq |O(t)|+\frac{k}{2}(j-1).
\end{align*}
In particular,
\begin{equation}
 j\leq\frac{2|O(t)|}{k}.
 \label{ha:component-count}
\end{equation}
This proves the decomposition estimate directly, including when $H$
is disconnected and without assuming that $k/4$ is an integer.

Let $U$ be the centers, identified with vertices, of the balls still
in $\operatorname{Int}_t$, and set
\[
 \begin{split}
 U_i&=\{u\in U:\{u,v\}\in O'(t)\text{ for some }v\in S_i\},\\
 V_i&=\{v\in S_i:\{u,v\}\in O'(t)\text{ for some }u\in U\}.
 \end{split}
\]
The definition of an assigned bag ensures $\bigcup_iU_i=U$.
By the loop condition and the $12C_3/k$ assignment bound,
\[
 |U|\geq\frac12|\operatorname{Bag}_t|
       \geq\frac{6C_3|O(t)|}{k}.
\]
Choose $i$ with $|U_i|$ maximal.  Even if the sets $U_i$ overlap,
their union bound and \eqref{ha:component-count} imply
\begin{equation}
 |U_i|\geq\frac{|U|}{j}\geq3C_3.
 \label{ha:many-neighbors}
\end{equation}
Choose a nonavoiding record $B$ with center $x$ and numerical radii
$r_1<r_2$, containing in its interior a ball with center $u_0\in U_i$.
Such a record exists because
every ball still under consideration is interior and operation 1 is
unavailable.  Write $\rho(z)=\|X_z-x\|_1$ and put
\begin{equation}
 a=\max\{r_1,\min_{v\in V_i}\rho(v)\},\qquad
 b=\min\{r_2,\max_{v\in V_i}\rho(v)\},\qquad L=b-a.
 \label{ha:radial-endpoints}
\end{equation}
The neighbor of $u_0$ in $V_i$ has radius strictly between
$r_1+(C_3-1)\delta$ and $r_2-(C_3-1)\delta$.  In particular $a\leq b$.
Define
\[
 \mathcal A=
 \{B(X_u,\delta)\in\operatorname{Int}_t:
                a-\delta<\rho(u)<b+\delta\},
 \qquad U_{\mathcal A}=\{u:B(X_u,\delta)\in\mathcal A\},
 \qquad m=|\mathcal A|.
\]
The set $\mathcal A$ is nonempty, since it contains the ball centered
at $u_0$.

Every ball $A\in\mathcal A$ is in the interior of $B$.  To check this
carefully, the radial inequalities defining $\mathcal A$ imply that
$A$ meets the support of $B$: its center is less than $\delta$ in
radial distance from $[r_1,r_2]$.  This assertion also holds for supports
restricted to the cube.  To decrease radius, move towards the binary
center $x$; to increase it, move towards the opposite cube vertex.
Along either segment, the change of radius equals the distance moved.
Since $B$ is nonempty, one can therefore reach a point of $B$ within
distance strictly less than $\delta$.  All balls of $\operatorname{Bag}_t$
have the same labels, so $C(A)$ intersects $C(B)$.  Also $A$ is
interior to some existing record $B'$.  If $B'\ne B$, then $B$ meets
$B'$ at a point of $A$, since the interior margin implies $A\subseteq
B'$.  Lemma~\ref{ha:unique-interior} contradicts
$C(A)\cap C(B)\ne\varnothing$.  Therefore $B'=B$, as claimed.

We need one more radial estimate:
\begin{equation}
 U_i\not\subseteq U_{\mathcal A}
       \quad\Longrightarrow\quad L\geq(C_3-1)\delta.
 \label{ha:radial-range}
\end{equation}
If every $v\in V_i$ had $r_1<\rho(v)<r_2$, then all their radii would
lie in $[a,b]$.  Each $u\in U_i$ has a neighbor $v\in V_i$ with
$|\rho(u)-\rho(v)|<\delta$, so it would lie in $U_{\mathcal A}$.
Hence the premise of \eqref{ha:radial-range} supplies some $v\in V_i$
with $\rho(v)\geq r_2$ or $\rho(v)\leq r_1$.  Let $w\in V_i$ be a
neighbor of the chosen interior center $u_0$.
If $\rho(v)\geq r_2$, then $b=r_2$, while
\[
 a\leq\rho(w)<\rho(u_0)+\delta
                  <r_2-(C_3-1)\delta.
\]
If $\rho(v)\leq r_1$, then $a=r_1$, while
\[
 b\geq\rho(w)>\rho(u_0)-\delta
                  >r_1+(C_3-1)\delta.
\]
These prove both cases of \eqref{ha:radial-range}.

\paragraph{\textbf{Operation 2: a dense group in the radial interval.}}
Suppose $m\delta>3L$.  Replace $B$ by
\[
 B_1=B(x,r_1\Vert a-2\delta),\qquad
 B_2=B(x,b+2\delta\Vert r_2),
\]
with its labels, and insert all $m$ original balls of $\mathcal A$,
with their original labels.  Treat these $m$ balls.  Their supports are
disjoint from $B_1,B_2$ because their centers have radii strictly
between $a-\delta$ and $b+\delta$.  They are pairwise disjoint as
members of the same family.  We proved that each lies in the interior
of $B$, so Lemma~\ref{ha:unique-interior} makes each compatible with
every old record other than $B$.  Their individual conditions follow
from insertability.  Residual shells inherit validity, proving that
the whole replacement is valid.

The number treated is large.  If $U_i\subseteq U_{\mathcal A}$, then
\eqref{ha:many-neighbors} gives $m\geq3C_3$.  Otherwise,
\eqref{ha:radial-range} and $m\delta>3L$ give
\begin{equation}
 m\geq3(C_3-1).
 \label{ha:dense-count}
\end{equation}
Allocate to $B$ all the potential of $B_1,B_2$ and three quarters of
the potential of each newly inserted ball.  Writing $R=r_2-r_1$,
the old potential is $R-C_4\delta$; it is positive because $B$
contains an interior ball.  Even if a residual shell is empty, its
potential is at least its formal width minus $C_4\delta$.  Therefore
\begin{align*}
 &\operatorname{tok}_\ell(B_1)+\operatorname{tok}_\ell(B_2)
       +\frac34\sum_{A\in\mathcal A}\operatorname{tok}_\ell(A)\\
 &\quad\geq R-L-4\delta-2C_4\delta+\frac58m\delta\\
 &\quad\geq\operatorname{tok}_\ell(B)
                -(C_4+4)\delta+\frac7{24}m\delta\\
 &\quad\geq\operatorname{tok}_\ell(B)
                -(C_4+4)\delta+\frac78(C_3-1)\delta\\
 &\quad\geq\operatorname{tok}_\ell(B).
\end{align*}
The successive estimates use \eqref{ha:fresh-potential},
$L<m\delta/3$, \eqref{ha:dense-count}, and
\eqref{ha:constant-consequences}.  Each inserted ball retains at least
$(1/4)(5/6)\delta=5\delta/24\geq\theta\delta$ to give to $t$.
Thus this operation supplies at least $\theta m\delta$ new credit.

\paragraph{\textbf{Operation 3: a sparse group in the radial interval.}}
Suppose $m\delta\leq3L$.  We first record that
\begin{equation}
 L\geq(C_3-1)\delta.
 \label{ha:sparse-width}
\end{equation}
If $U_i\not\subseteq U_{\mathcal A}$ this follows from
\eqref{ha:radial-range}; otherwise
$L\geq m\delta/3\geq|U_i|\delta/3\geq C_3\delta$.
In particular, all the central shells defined next have positive width.

Replace $B$ by the two residual shells
\[
 B_1=B(x,r_1\Vert a),\qquad B_2=B(x,b\Vert r_2),
\]
with its labels.  Insert two additional shells
\[
 B_3=B(x,a+\delta\Vert b-\delta),\qquad
 B_4=B(x,a\Vert b).
\]
The shell $B_3$ is nonavoiding and $B_4$ is avoiding.  Give them labels
\begin{equation}
 \begin{array}{c|ccc}
    &t(\,\cdot\,)&t_d(\,\cdot\,)&t_P(\,\cdot\,)\\ \hline
 B_3&t&\text{none}&
             \{v\in t_P(B):v\text{ is a proper descendant of }t\}\\
 B_4&t(B)&t&t_P(B).
 \end{array}
 \label{ha:central-labels}
\end{equation}
Treat all the balls of $\mathcal A$, removing them from
$\operatorname{Int}_t$ without inserting them.

We verify every new labeling condition.  Lemma~\ref{ha:ancestor}
implies that $t(B)$ is a weak ancestor of $t$ and that $t\in C(B)$.
In fact $t(B)\ne t$.  An original ball with owner $t$ cannot meet an
unprocessed ball of the same family, since that family is disjoint.
A synthetic shell with owner $t$ could only have been created by an
earlier occurrence of this operation while processing this very bag:
the sets $T_i$ are disjoint, and new nonavoiding owners are assigned
only to the node being processed.  At that earlier occurrence every
remaining interior ball meeting the new central shell was in its
expanded radial interval and was treated.  Subsequent changes only
restrict that shell, and Lemma~\ref{ha:border-monotone} prohibits the
return of a border ball to the interior set.  Hence neither that shell
nor any of its residual pieces can now contain an unprocessed
interior ball.  This proves $t(B)\ne t$.  Thus $t$ is a proper
descendant of $t(B)$, as required for $B_4$ to be avoiding.

Using $\mathcal T[t]$ for the subtree rooted at $t$, the labels give
\begin{equation}
 C(B_3)=C(B)\cap\mathcal T[t],
 \qquad C(B_4)=C(B)\setminus\mathcal T[t].
 \label{ha:conflict-partition}
\end{equation}
Indeed, no member of $t_P(B)$ is a weak ancestor of $t$, since
$t\in C(B)$; every remaining excluded subtree either lies below $t$
or is disjoint from its subtree.  Both new conflict domains have the
prescribed connected form and are subsets of $C(B)$.  All four shells
are subsets of $B$.  The pairwise validity with old records follows
from Lemma~\ref{ha:heredity}; the residual shells are disjoint from
the central shells, and \eqref{ha:conflict-partition} proves validity
of the overlapping pair $B_3,B_4$.  Finally, both new excluded-node
sets are subsets of $t_P(B)$ and both new outer radii are at most
$r_2$.  Thus the excluded-boundary-endpoint condition is preserved.

It remains to prove the crossing-path condition, which does not follow
merely from shrinking the geometric support when labels change.
Choose vertices $v_1,v_2\in V_i$ minimizing and maximizing $\rho$, and
choose their respective neighbors $u_1,u_2\in U_i$ along edges of
$O'(t)$.  Clipping in \eqref{ha:radial-endpoints} gives
\begin{equation}
 \rho(v_1)\leq a,\quad \rho(v_2)\geq b,
 \qquad \rho(u_1)<a+\delta,\quad \rho(u_2)>b-\delta.
 \label{ha:radial-witnesses}
\end{equation}
By \eqref{ha:sparse-width}, $a+\delta<b\leq r_2$.  Hence
$\rho(u_1)<r_2$ and $\rho(v_1)<r_2$; these strict inequalities will
exclude membership in the forbidden subtrees.

For $B_3$, the $k$-edge-connectivity of $G[V(t)]$ and
Lemma~\ref{ha:edge-paths} supply $k$ edge-disjoint $u_1$--$u_2$ paths.
The radial inequalities show that
each crosses $B_3$.  To see that the initial vertex is in $C(B_3)$,
fix $z\in t_P(B_3)$.  If $u_1\in V(z)$, the edge
$\{u_1,v_1\}\in O'(t)$ would belong to $P(z)$ because $z$ is below
$t$ and $v_1\notin V(t)$.  This contradicts the endpoint condition of
$B$, since $z\in t_P(B)$ but $\rho(u_1)<r_2$.  Thus $u_1$ avoids
every excluded subtree and lies in $C(B_3)$.  These paths and this
initial vertex prove the required crossing condition for $B_3$.

For $B_4$, every cut separating $v_1$ from $v_2$ in $G[S_i]$ has
at least $\lceil k/4\rceil$ edges.  Lemma~\ref{ha:edge-paths}
therefore supplies at least $k/4$ edge-disjoint $v_1$--$v_2$ paths.
Since $t(B)$ is a proper
ancestor of $t$, $V(t^*)\subseteq V(t(B))$.  The paths therefore lie
in $V(t(B))\setminus V(t)$ and cross $B_4$.  Fix $z\in t_P(B_4)$.
The edge $\{u_1,v_1\}$ cannot belong to $P(z)$ because
$\rho(v_1)<r_2$.  If $v_1\in V(z)$, it follows that $u_1\in V(z)$
also.  A hierarchy set containing both $u_1\in V(t)$ and
$v_1\notin V(t)$ must be a proper ancestor of $t$, by laminarity.
This would make $z$ a weak ancestor of $t$, contrary to $t\in C(B)$.
Hence $v_1$ lies in none of the excluded subtrees; also
$v_1\notin V(t)$, so $v_1\in C(B_4)$.  This verifies the crossing
condition for $B_4$ and completes the proof of validity.

For later use, every still-unprocessed interior ball is disjoint from
$B_3$ and $B_4$: its center is outside the expanded interval
$(a-\delta,b+\delta)$.  Such a ball can remain interior to $B_1$ or
$B_2$, which is harmless.  This establishes the property of new
central shells used above to prove that $t(B)$ is a proper ancestor.

We now distribute the new potentials.  The shell $B_3$, although
created in the current phase, is not an original current-family ball;
it uses the second line of \eqref{ha:potential}.  Give $B$ all the
potential of $B_1,B_2,B_3$ and $(2C_4+2)\delta$ from $B_4$.
The sum of the three formal nonavoiding widths is $R-2\delta$,
where $R=r_2-r_1$.  Consequently
\begin{align*}
 \sum_{j=1}^3\operatorname{tok}_\ell(B_j)+(2C_4+2)\delta
 &\geq R-2\delta-3C_4\delta+(2C_4+2)\delta\\
 &=R-C_4\delta=\operatorname{tok}_\ell(B).
\end{align*}
The following estimate proves both that $B_4$ has the amount just
promised and that it can additionally pay for all $m$ treated balls:
\begin{align*}
 \operatorname{tok}_\ell(B_4)-(2C_4+2)\delta
 &=\frac{L-C_4\delta-4(C_4+1)(C_4+2)\delta}{2(C_4+2)}\\
 &\geq\frac{L-\{C_4+4(C_4+2)^2\}\delta}{2(C_4+2)}\\
 &\geq\frac{L-(C_3-1)\delta/2}{2(C_4+2)}\\
 &\geq\frac{L}{4(C_4+2)}
 \geq\frac{m\delta}{12(C_4+2)}
 \geq\theta m\delta.
\end{align*}
The first equality is valid because \eqref{ha:sparse-width} makes
$L-C_4\delta$ positive.  The other steps use
$C_4+1\leq C_4+2$, \eqref{ha:constant-consequences},
\eqref{ha:sparse-width}, the case assumption $m\delta\leq3L$, and
$C_3\geq2C_4(C_4+2)$, respectively.  Give the remaining amount to $t$.
This completes the allocation for operation 3.

Each operation treats at least one interior ball, preserves all old
allocations, and gives $t$ at least $\theta\delta$ for each ball it
treats.  Border transfers change neither this count nor any allocation.
Induction on the number of operations proves the stated bound for
$t$, and induction on the bags already processed in the phase proves
the bounds for the other current-phase nodes.  The loop terminates
after at most $|\operatorname{Bag}_t|$ operations; its negated stopping
condition gives the asserted strict half-bag bound.  This proves the
lemma.  The first phase starts with an empty arrangement and no old
allocations; the transfer of these allocations between phases is
verified in the next subsection.
\end{proof}
